\section*{Chapitre \ref{ch:b1:magnetostatics} --- Magnétostatique~: champ, forces et dipôles}

\begin{solution}{exo:b1:magnetostatics:1}
$B = \mu_0I/2\pi r = 2 \times 10^{-7} \times 10/0.01 = \qty{2e-4}{T}$. Pour \qty{1}{T}~:
$I = 2\pi rB/\mu_0 = \qty{5e4}{A}$. La valeur terrestre est atteinte en $r = 2 \times 10^{-7} \times
100/5 \times 10^{-5} = \qty{0.4}{m}$.
\end{solution}

\begin{solution}{exo:b1:magnetostatics:2}
Centre~: $\mu_0NI/2R = 4\pi \times 10^{-7} \times 200/0.1 = \qty{2.5}{mT}$. En $z = R$~:
$\times(R^2/2R^2)^{3/2} = 2^{-3/2}$, soit \qty{0.89}{mT}. En $z = 10R$~: $\times(1/101)^{3/2}
= 10^{-3}$, soit \qty{2.5}{\micro T}, déjà la loi dipolaire en $1/z^3$.
\end{solution}

\begin{solution}{exo:b1:magnetostatics:3}
$B = \mu_0nI = 4\pi \times 10^{-7} \times 1000 \times 5 = \qty{6.3}{mT}$~; au bout,
la moitié~: \qty{3.1}{mT}. Avec du fil de \qty{0.5}{mm}~: $2000$ spires par mètre,
donc \qty{12.6}{mT} sous \qty{5}{A}.
\end{solution}

\begin{solution}{exo:b1:magnetostatics:4}
$F = ILB = 10 \times 1 \times 0.1 = \qty{1.0}{N}$~; à \qty{30}{\degree}, il faut $\sin30^\circ$~:
\qty{0.5}{N}. Moteur~: $100 \times \qty{1}{N} \times \qty{0.05}{m} = \qty{5}{N.m}$ ---
un petit moteur industriel.
\end{solution}

\begin{solution}{exo:b1:magnetostatics:5}
Avec $z = d\tan\alpha$~: $B = \dfrac{\mu_0I}{4\pi}\displaystyle\int\dfrac{d\,\dd z}{(d^2 + z^2)^{3/2}}
= \dfrac{\mu_0I}{4\pi d}\displaystyle\int_{\alpha_1}^{\alpha_2}\cos\alpha\,\dd\alpha = \dfrac{\mu_0I}{4\pi d}
(\sin\alpha_2 - \sin\alpha_1)$~; fil infini, $\alpha_{1,2} = \mp\pi/2$~: $\mu_0I/2\pi d$.
Carré~: chaque côté, à $d = a/2$ et vu sous $\pm45^\circ$, donne $\mu_0I\sqrt2/2\pi a$~;
pour les quatre côtés~: $B = 2\sqrt2\,\mu_0I/\pi a = 0.90\,\mu_0I/a$ (un cercle de même
«~rayon~» $a/2$ donnerait $\mu_0I/a$).
\end{solution}

\begin{solution}{exo:b1:magnetostatics:6}
$B(z) = f(z - R/2) + f(z + R/2)$ avec $f(u) = \mu_0NIR^2/2(R^2 + u^2)^{3/2}$.
Au milieu~: $2f(R/2) = \mu_0NIR^2/(5R^2/4)^{3/2} = (4/5)^{3/2}\mu_0NI/R$. $f'$ est
impaire, donc $B'(0) = f'(-R/2) + f'(R/2) = 0$~; $f''(u) \propto (4u^2 - R^2)(R^2 +
u^2)^{-7/2}$ s'annule en $u = R/2$, donc $B''(0) = 2f''(R/2) = 0$~: le champ est
uniforme au troisième ordre. Application~: $0.716 \times 4\pi \times 10^{-7} \times 100/0.15 =
\qty{0.60}{mT}$.
\end{solution}

\begin{solution}{exo:b1:magnetostatics:7}
Surface~: $\mu_0I/2\pi a = 2 \times 10^{-7} \times 100/0.002 = \qty{10}{mT}$~; à
\qty{1}{mm}, à l'intérieur, $\times r/a$~: \qty{5}{mT}~; à \qty{1}{cm}~: \qty{2}{mT}. $B$ croît
linéairement jusqu'à la surface, puis décroît en $1/r$. Coaxial~: à l'extérieur
de la \omterm{def:b1:rays-reflection-refraction:fiber}{gaine} le courant enlacé est nul, donc $B = 0$.
\end{solution}

\begin{solution}{exo:b1:magnetostatics:8}
$2 \times 10^{-7}$~N/m. Jeux de barres~: $2 \times 10^{-7} \times 10^8/0.1 = \qty{200}{N/m}$~;
sous \qty{100}{kA}~: $\qty{2e4}{N/m}$ --- deux tonnes par mètre. Les supports
isolants de l'appareillage sont dimensionnés pour le courant de défaut, non
pour celui de service.
\end{solution}

\begin{solution}{exo:b1:magnetostatics:9}
$\Gamma = mB_h = \qty{2e-6}{N.m}$~; $W = 2mB_h = \qty{4e-6}{J}$~; $J\ddot\theta =
-mB_h\theta$~: $T = 2\pi\sqrt{J/mB_h} = 2\pi\sqrt{10^{-7}/2 \times 10^{-6}} =
\qty{1.4}{s}$. On chronomètre les oscillations~: $B_h = 4\pi^2J/mT^2$ une fois $m$ et $J$
connus (Gauss mesurait $m$ séparément, par la déviation qu'il imprime à une
seconde aiguille).
\end{solution}

\begin{solution}{exo:b1:magnetostatics:10}
(a) $\mu_0IR^2/2z^3 = \mu_0(I\pi R^2)/2\pi z^3 = \mu_0m/2\pi z^3 = (\mu_0/4\pi)\,2m/z^3$~:
c'est le $2p/4\pi\varepsilon_0z^3$ électrique avec $1/\varepsilon_0 \to \mu_0$. (b) $m = 2\pi R_E^3B/
\mu_0 = 2\pi \times 2.58 \times 10^{20} \times 6 \times 10^{-5}/1.26 \times 10^{-6} = \qty{7.7e22}{A.m^2}$.
(c) $I = m/\pi R^2 = 7.7 \times 10^{22}/(\pi \times 9 \times 10^{12}) = \qty{2.7e9}{A}$. (d)
Équateur ($\theta = 90^\circ$)~: $B_\theta = \mu_0m/4\pi R_E^3 = \qty{3e-5}{T}$, la moitié
de la valeur polaire. En $\theta = 45^\circ$~: $\tan(\text{inclinaison}) = B_r/B_\theta = 2\cot\theta = 2$,
inclinaison $= \qty{63}{\degree}$ --- aux latitudes moyennes, le champ plonge
fortement dans le sol.
\end{solution}

\begin{solution}{exo:b1:magnetostatics:11}
$\mu_0nI = 4\pi \times 10^{-7} \times 2000 \times 3 = \qty{7.5}{mT}$. Au centre~: $\tan\alpha_1 =
R/(L/2) = 0.2$, $\cos\alpha_1 = -\cos\alpha_2 = 0.981$~: $B = 0.981\,\mu_0nI =
\qty{7.4}{mT}$ ($98\%$). Au bout~: $\cos\alpha_1 = 0$, $\tan\alpha_2 = R/L$, $\cos\alpha_2 =
-0.995$~: $B = 0.497\,\mu_0nI = \qty{3.7}{mT}$ ($50\%$). À l'extérieur, en $x$ au-delà de
l'extrémité, les deux bouts étant du même côté~: $B = \tfrac12\mu_0nI(\cos\alpha_1 - \cos\alpha_2)
\approx \tfrac12\mu_0nI\,\tfrac{R^2}{2}[1/x^2 - 1/(x + L)^2]$~; $1\%$ de la valeur
centrale vers $x \approx \qty{10}{cm}$, une demi-longueur plus loin. $B(z)$~: un plateau
plat sur presque toute la longueur, qui tombe à la moitié aux extrémités et à
quelques pour cent une longueur plus loin.
\end{solution}

\begin{solution}{exo:b1:magnetostatics:12}
(a) Le champ propre des enroulements saute de $0$ à $B$ à leur traversée~; le
champ qui agit sur eux est la moyenne, $B/2$ (une nappe ne peut pas se pousser
elle-même). Force sur un élément d'enroulement de longueur $\dd l$ et de largeur $\dd z$
parcouru par $K\,\dd z$~: $K\,\dd z\,\dd l\,B/2$, radiale vers l'extérieur (à vérifier avec $I\,\dd\vect l
\wedge\vect B$)~: pression $p = KB/2 = B^2/2\mu_0$ puisque $K = nI = B/\mu_0$.
(b) \qty{1}{T}~: \qty{4e5}{Pa} = \qty{4}{bar}~; \qty{10}{T}~: \qty{400}{bar}~; \qty{45}{T}~:
\qty{8e8}{Pa} = \qty{8000}{bar}, proche du \qty{1}{GPa} de l'acier. (c) Les
enroulements doivent tenir une pression qui croît comme $B^2$~: au-delà de
\qty{40}{T} environ, aucun matériau ni aucun refroidissement n'y suffit --- la
limite des aimants à champ continu est mécanique (et, pour les
supraconducteurs, leur champ critique).
\end{solution}

\begin{solution}{pb:b1:magnetostatics:1}
\textbf{1.} Tout plan contenant l'axe est un \omterm{prop:b1:electrostatics-gauss:symmetry}{plan de symétrie} des courants~:
$\vect B$ lui est perpendiculaire, donc orthoradial~; par invariance le long
de l'axe et autour de lui~: $\vect B = B(r)\,\vect e_\theta$.

\textbf{2.} Cercle de rayon $r$~: $2\pi rB = \mu_0I_{\text{enlacé}}$. Dans l'âme~:
$I_{\text{enlacé}} = Ir^2/a^2$, $B = \mu_0Ir/2\pi a^2$~; entre les conducteurs~: $\mu_0I/2\pi r$~;
au-dehors ($r > c$)~: $I - I = 0$, donc $B = 0$.

\textbf{3.} $B(a) = 2 \times 10^{-7}/5 \times 10^{-4} = \qty{4.0e-4}{T}$~; $B(b) = 2 \times
10^{-7}/2.5 \times 10^{-3} = \qty{8.0e-5}{T}$.

\textbf{4.} À l'intérieur du tube, aucun courant n'est enlacé~: $B = 0$ au lieu
de la croissance linéaire~; partout ailleurs rien ne change --- le champ à
l'extérieur d'un courant cylindrique ne dépend que du courant total.

\textbf{5.} $I_{\text{enlacé}} = I - I(r^2 - b^2)/(c^2 - b^2) = I(c^2 - r^2)/(c^2 -
b^2)$~: $B = \dfrac{\mu_0I}{2\pi r}\,\dfrac{c^2 - r^2}{c^2 - b^2}$, qui vaut $\mu_0I/2\pi b$ en
$b$ et $0$ en $c$~: c'est continu. Tracé~: croissance linéaire jusqu'à \qty{4e-4}{T} à
\qty{0.5}{mm}, décroissance en $1/r$ jusqu'à \qty{8e-5}{T} à \qty{2.5}{mm}, chute à zéro à
\qty{3}{mm}, nul au-delà.

\textbf{6.} Le courant de la \omterm{def:b1:rays-reflection-refraction:fiber}{gaine} est de sens opposé à celui de l'âme~: répulsion,
donc une pression vers l'extérieur. Courant surfacique $K = I/2\pi b = \qty{64}{A/m}$~; champ moyen
$\tfrac12(8 \times 10^{-5} + 0) = \qty{4e-5}{T}$~: $p = KB = \qty{2.5e-3}{Pa}$. Sous
\qty{10}{kA}~: $\times10^8$, soit \qty{2.5e5}{Pa} = \qty{2.5}{bar} --- les câbles de
puissance impulsionnelle sont construits pour l'encaisser.

\textbf{7.} Aucun champ au-dehors~: le câble ne rayonne pas et ne capte pas les
perturbations magnétiques de ses voisins. Pour une paire de fils, les deux
champs se compensent presque, laissant $\approx \mu_0Id/2\pi D^2 = 2 \times 10^{-7} \times 2 \times
10^{-3}/10^{-2} = \qty{4e-8}{T}$ à \qty{10}{cm} --- c'est peu, mais le coaxial
donne zéro.

\textbf{8.} $D = \sqrt{h^2 + d^2/4} = \qty{10.0}{m}$~: chacun donne $\mu_0I/2\pi D =
\qty{1.0e-5}{T}$, perpendiculaire à la droite qui joint le conducteur au point.
Les courants étant opposés, les composantes parallèles à la droite qui joint
les deux conducteurs se compensent et les verticales s'ajoutent~: $B = 2 \times 10^{-5} \times
(d/2)/D = \qty{1.0e-6}{T}$, vertical. Cinquante fois sous le champ terrestre,
cent fois sous la valeur repère.

\textbf{9.} Pour $D \gg d$, $B \approx \dfrac{\mu_0I}{2\pi}\Bigl(\dfrac1{D_1} - \dfrac1{D_2}\Bigr)
\sim \dfrac{\mu_0I}{2\pi}\dfrac{d}{D^2}$~: les deux champs en $1/D$ se compensent
au premier ordre et seule survit leur différence, d'ordre relatif $d/D$ --- un
«~dipôle à deux fils~».

\textbf{10.} Son propre conducteur à \qty{1}{m}~: \qty{1.0e-4}{T} horizontal~; l'autre
à \qty{1.41}{m}~: \qty{7.1e-5}{T} à \qty{45}{\degree}, en sens contraire~: résultante
$(5 \times 10^{-5}, 5 \times 10^{-5})$, soit \qty{7e-5}{T} --- \qty{70}{\micro T}, sous la
valeur repère du public et très loin sous la valeur professionnelle.

\textbf{11.} $F/\ell = \mu_0I^2/2\pi d = 2 \times 10^{-7} \times 2.5 \times 10^5 =
\qty{0.05}{N/m}$, répulsive~; sous \qty{20}{kA}~: \qty{80}{N/m}.

\textbf{12.} \qty{15}{N} en régime normal~; \qty{24}{kN} en court-circuit, huit fois
le poids de la portée (\qty{2.9}{kN})~: les conducteurs sont écartés violemment,
balancent et peuvent s'entrechoquer en retombant --- d'où les entretoises et
des disjoncteurs qui ouvrent en quelques périodes.

\textbf{13.} $I\cos\omega t + I\cos(\omega t - 2\pi/3) + I\cos(\omega t + 2\pi/3) =
I\cos\omega t\,(1 + 2\cos2\pi/3) = 0$~: les trois \omterm{def:b1:sinusoidal-impedance:complex}{vecteurs de Fresnel} ont une
somme nulle. Les champs en $1/D$ se compensent comme pour le coaxial~; le champ
lointain décroît en $1/D^2$ ou plus vite.

\textbf{14.} Le champ de la ligne y est vertical~; une boussole ne répond qu'à
la composante horizontale~: aucune déviation (l'inclinaison change de
$10^{-6}/(2 \times 10^{-5})$, soit trois degrés, que la boussole ne montre pas).

\textbf{15.} Côtés $a$ (parallèles à l'axe)~: le $\vect B$ radial leur est
perpendiculaire, d'où une force $NIaB$ sur chacun, perpendiculaire au plan du
cadre et opposée sur les deux côtés~: un couple. Côtés $b$~: $\dd\vect l$ y est
radial, donc parallèle à $\vect B$~: aucune force.

\textbf{16.} $\Gamma = 2 \times NIaB \times b/2 = NIabB = NISB = 0.024\,I$ (en N\,m, $I$
en A). Le champ radial est toujours dans le plan du cadre et perpendiculaire
aux côtés $a$~: les forces sont toujours normales au plan et leur \omterm{def:b1:angular-momentum:moment}{bras de
levier} $b/2$ ne change jamais.

\textbf{17.} $C\theta = NISB$~: $\theta = 2400\,I$ rad, soit \qty{2.4}{rad/mA}~; pleine
échelle à \qty{1.2}{rad} pour $I = \qty{0.50}{mA}$.

\textbf{18.} Longueur $200 \times 2(0.02 + 0.03) = \qty{20}{m}$, section $\pi(5 \times
10^{-5})^2 = \qty{7.9e-9}{m^2}$~: $R = 1.7 \times 10^{-8} \times 20/7.9 \times 10^{-9} =
\qty{43}{\ohm}$~; $U = 43 \times 5 \times 10^{-4} = \qty{22}{mV}$~; $P = UI = \qty{11}{\micro W}$.

\textbf{19.} $J\ddot\theta = -C\theta + NISB\,I$~: $\omega_0 = \sqrt{C/J} = \sqrt{10^{-5}/6.7
\times 10^{-8}} = \qty{12}{rad/s}$, $T = \qty{0.51}{s}$.

\textbf{20.} $J\ddot\theta + \alpha\dot\theta + C\theta = NISB\,I$~: critique pour
$\alpha = 2\sqrt{JC} = 2\sqrt{6.7 \times 10^{-13}} = \qty{1.6e-6}{N.m.s}$~; l'aiguille
atteint alors sa valeur le plus vite possible sans dépassement.

\textbf{21.} Voltmètre~: $R_s = 10/5 \times 10^{-4} - 43 \approx \qty{20}{k\ohm}$.
Ampèremètre~: le shunt porte \qty{0.9995}{A} sous \qty{22}{mV}, d'où $R_{\text{sh}} =
\qty{22}{m\ohm}$.

\textbf{22.} $m = NIS = 200 \times 5 \times 10^{-4} \times 6 \times 10^{-4} = \qty{6e-5}{A.m^2}$~;
$-mB = \qty{-1.2e-5}{J}$~; ressort $\tfrac12C\theta^2 = \tfrac12 \times 10^{-5} \times 1.44 =
\qty{7.2e-6}{J}$~: le même ordre --- le ressort emmagasine ce que le couple
magnétique a fourni pendant la rotation du cadre.

\textbf{23.} Dans un champ uniforme, le moment vaut $mB\cos\theta$ (maximal quand
le plan contient $\vect B$, nul quand la normale lui est parallèle)~: $C\theta =
NISB\,I\cos\theta$ --- linéaire seulement aux petits $\theta$, et tassé en haut
de l'échelle.

\textbf{24.} \qty{50}{Hz} contre $\omega_0/2\pi = \qty{2}{Hz}$~: le cadre ne peut
pas suivre~; il reste à la moyenne du courant, c'est-à-dire zéro, avec un
frémissement d'\omterm{def:b1:wave-propagation:signal}{amplitude} $\approx \theta_{\text{statique}}(\omega_0/\omega)^2 = 1.2 \times (12/314)^2 \approx
\qty{2e-3}{rad}$, invisible. Un appareil à cadre mobile indique la moyenne ---
pour l'alternatif, il lui faut un redresseur.

\textbf{25.} \qty{4e-4}{T} au cœur d'un coaxial de \qty{1}{A} et rien au-dehors
(\omterm{thm:b1:magnetostatics:ampere}{théorème d'Ampère})~; \qty{1}{\micro T} sous une ligne de \qty{500}{A} et
\qty{0.05}{N/m} entre ses conducteurs (Biot et Savart, puis Laplace)~;
\qty{2.4}{rad/mA} pour l'appareil (couple de Laplace contre un ressort).
\end{solution}
